% SPDX-FileCopyrightText: Copyright (c) 2025-2026 Objectionary.com
% SPDX-License-Identifier: MIT

Java 8 introduced lambda expressions and streams in 2014.
Streams make code more expressive compared to traditional imperative loops~\citep{khatchadourian2020study}.
Although they resemble functions from functional programming, they are objects, not functions.
The Stream API consists of regular Java classes, while the lambda expressions passed to stream operations are compiled to \ff{invokedynamic} instructions and corresponding static methods.
At runtime, when the JVM encounters an \ff{invokedynamic} instruction, it creates an instance of a class that implements the target functional interface.
All calls to the lambda expression then become virtual calls to that created instance.

Streams in Java can be slow due to their abstraction overhead~\citep{kiselyov2017,khatchadourian2020study}.
Each \ff{map} and \ff{filter} operation creates a new stream with its own infrastructure.
For a simple sum of five squares, the stream pipeline invokes the lambda function five times, plus additional overhead from stream machinery:
\begin{ffcode}
var NN = {1, 2, 3, 4, 5};
var s = IntStream.of(NN).map((*@\highlight{x -> x * x}@*)).sum();
\end{ffcode}
This algorithm can be rewritten as an imperative loop that doesn't use virtual calls:
\begin{ffcode}
var s = 0;
for (var i = 0; i < NN.length; i++)
  s += NN[i] * NN[i];
\end{ffcode}

Performance advantages of the imperative solution were obvious in earlier JVMs.
However, modern JVMs are different.

\begin{figure*}
\settwocapwidths{0.45}
\begin{subfigure}[t]{\firstsf}
\begin{ffcode}
public long loop() {
  long sum = 0L;
  (*@\highlight{for}@*) (int i=0; i<NN.length; ++i)
    if (NN[i] % 2 == 0)
      sum += NN[i] * NN[i];
  return sum; }
\end{ffcode}
\end{subfigure}
\hfill
\begin{subfigure}[t]{\secondsf}
\begin{ffcode}
public long map() {
  return IntStream.of(NN)
    .(*@\highlight{filter}@*)(x -> x % 2 == 0)
    .(*@\highlight{map}@*)(x -> x * x)
    .(*@\highlight{sum}@*)(); }
\end{ffcode}
\end{subfigure}
\twofigcaps
  {moller}
  {Imperative style.}
  {Functional style, using a stream pipeline.}
  {Two variants of computing sums of even squares suggested by \citet{moller2020eliminating} to demonstrate that the usage of the Stream API degrades performance.}
\end{figure*}

\begin{table}
\setlength{\tabcolsep}{.3em}
\figcap{%
  The ``Ratio'' column shows stream time divided by loop time.
  It compares \cref{fig:moller-left} with \cref{fig:moller-right} for each JVM version.
  The numbers demonstrate that the performance of streams was improved since earlier JVM versions.}
\intabular{moller-jvms}
\end{table}

Modern JVMs handle streams better than their earlier versions.
We reproduced the experiment of \citet{moller2020eliminating}.
\Cref{fig:moller} shows two code snippets they tested with Java~13 using the OpenJDK Server~VM (build 13+33).
They observed a ``60\% performance degradation'' for the stream approach (\cref{fig:moller-right}) compared to the loop (\cref{fig:moller-left}).
\cref{tab:moller-jvms} shows the data we collected: the degradation is present only in versions older than Java~11.
Moreover, some of the latest JVMs calculate the sum of even numbers faster when using a stream of integers.

In this and all other experiments, we used JMH\footnote{\url{https://github.com/openjdk/jmh}} with
  \input{_env/warmups.tex} warming-up and \input{_env/iterations.tex} measuring iterations.
We used
  a bare metal
  \input{_env/os.tex}
  at
  \input{_env/ghz.tex}~GHz
  \input{_env/arch.tex}
  machine
  with \input{_env/cores.tex}~cores
  and \input{_env/ram.tex}~GiB of memory.
The maximum measurement error observed in all experiments was \input{_env/mistake.tex}\%.
The size of the \ff{NN} array was \input{_env/numbers.tex} million.
The relative standard deviation (RSD), here and in all other tables in the paper, is a measure of relative variability.
We performed all experiments \input{_env/repeat.tex} times and calculated RSD from the standard deviation and mean.
In all our measurement experiments we used JVMs installable via SdkMan\footnote{\url{https://sdkman.io/}}.

Under the hood, modern HotSpot JVMs have dramatically improved their JIT compilers and runtime machinery compared to the Java~8 era.

First, their inlining is more aggressive.
In Java~8, the C2 compiler's default MaxInlineLevel was 9, limiting nested call inlining.
By JDK~11, the inlining budget and heuristics were refined.
They better handle functional patterns.
Small lambdas and short stream pipeline operations are inlined and optimized into efficient loops.

Second, they improved escape analysis and scalar replacement.
Intermediate objects created by the Stream API can now be optimized away.
If the JIT proves that they never escape the method boundary, it eliminates the allocation entirely through scalar replacement.
This transforms what used to be dozens of fleeting allocations and GC events into pure register operations.

By contrast, a manually written loop often presents less uniform control flow.
It hides opportunities for scalar replacement and may not be profiled as comprehensively by the compiler.
Short pipelines---once vilified for their ``object churn''---now run as fast as or faster than equivalent for-loops.
